\chapter{Discussion}

\section{Lectures principales}
Trois constats ressortent des campagnes disponibles:
\begin{itemize}
  \item l'\textbf{occlusion} est systematiquement competitive et souvent premiere dans les protocoles les plus contraints;
  \item les combinaisons lourdes d'augmentations ne garantissent pas de gain stable par rapport a une politique plus simple;
  \item la configuration du monitor d'early stopping (\texttt{val\_acc} vs \texttt{val\_loss}) influence notablement l'epoch retenue et les metriques finales.
\end{itemize}

\section{Interpretation}
L'occlusion aleatoire est pertinente pour ModelNet10 car elle force le reseau a reconnaitre des formes partielles.
Cela ameliore la robustesse a la disparition locale de points, ce qui est coherent avec des scenes reelles ou les acquisitions sont incompletes.
Le gain en test du notebook final (\(90.09\%\)) valide ce choix pour la version courante du projet.

\section{Limites actuelles}
\begin{itemize}
  \item les benchmarks robustes ne couvrent pas encore un vrai multi-seeds complet (plusieurs tables ont \texttt{num\_runs = 1});
  \item la comparaison Exercice 1 / Exercice 2 reste surtout architecturale dans cette version du rapport, faute de tableaux consolides homogènes sur la meme grille d'essais;
  \item l'evaluation est concentree sur ModelNet10, sans extension ModelNet40.
\end{itemize}

\section{Etat en cours}
Au moment de redaction, un fine-tuning \texttt{baseline + occlusion} en \texttt{seed 45} est en cours.
Il est prevu d'ajouter ses metriques finales (accuracy, loss, matrice de confusion) dans le tableau de comparaison principal des que le run est termine.
